% 太阳系（综述）
% license CCBYSA3
% type Wiki

本文根据 CC-BY-SA 协议转载翻译自维基百科\href{https://en.wikipedia.org/wiki/Solar_System}{相关文章}。

\begin{figure}[ht]
\centering
\includegraphics[width=8cm]{./figures/5a727e2db62d7a0e.png}
\caption{太阳、行星、卫星与矮行星\(^\text{[a]}\)（真实色彩，比例为尺寸等比，但距离不按比例）} \label{fig_Solar_1}
\end{figure}
太阳系（The Solar System）由太阳及其所环绕的天体组成。该名称来源于“Sōl”，即太阳的拉丁名称。大约在 46 亿年前，一个分子云中密度较高的区域发生坍缩，形成了太阳以及一个原行星盘（protoplanetary disc），环绕天体便由此汇聚而成。太阳核心内部的氢融合为氦，释放出能量，这些能量主要通过其外层光球（photosphere）向外辐射，从而在整个系统中形成一个由内而外逐渐降低的温度梯度。太阳系中超过 99.86\% 的质量集中在太阳内部。

在绕太阳运行的天体中，质量最大的为八大行星。按与太阳距离递增的顺序，最靠近太阳的四颗类地行星（terrestrial planets）分别是—— 水星（Mercury）、金星（Venus）、地球（Earth）和火星（Mars）。它们构成了太阳系的内行星系统（inner Solar System）。其中，地球与火星是太阳系中仅有的两颗位于太阳宜居带（habitable zone）内的行星——在此范围内，行星表面可存在液态水。

在距离太阳约 5 个天文单位（AU）之外、越过“霜线”（frost line）的位置，分布着四颗更为庞大的行星：两颗气态巨行星（gas giants）—— 木星（Jupiter）与土星（Saturn），以及两颗冰巨行星（ice giants）—— 天王星（Uranus）与 海王星（Neptune）。它们共同构成了太阳系的外行星系统（outer Solar System）。其中，木星与土星占据了太阳系除恒星外几乎 90\% 的全部质量。

太阳系中还存在着数量庞大但质量较小的天体。天文学家普遍认为，太阳系至少拥有九颗矮行星（dwarf planets），分别是：谷神星（Ceres）、奥库斯（Orcus）、冥王星（Pluto）、哈梅亚（Haumea）、夸奥尔（Quaoar）、鸟神星（Makemake）、共工（Gonggong）、阋神星（Eris）与赛德娜（Sedna）。在这些天体中，有六颗行星、七颗矮行星及若干其他天体拥有自然卫星（natural satellites），即通常所称的“月亮（moons）”。这些卫星的大小范围极广：从如地球之月那样接近矮行星尺度的巨型卫星，到仅有极小质量的“月粒（moonlets）”。此外，太阳系中还存在大量小型天体（small Solar System bodies），包括小行星（asteroids）、彗星（comets）、半人马天体（centaurs）、流星体（meteoroids）以及星际尘埃云（interplanetary dust clouds）。其中一部分天体分布在小行星带（asteroid belt）（位于火星与木星轨道之间）以及柯伊伯带（Kuiper belt）（位于海王星轨道之外）。

在这些天体之间，存在着充满尘埃与粒子的行星际介质（interplanetary medium）。太阳不断释放的带电粒子形成太阳风（solar wind），它向外扩散，构成一个被称为日球层（heliosphere）的巨大空间。在距离太阳约 70–90 个天文单位（AU）处，太阳风受到星际介质（interstellar medium）的阻挡，形成边界——日球顶（heliopause）。这一区域标志着太阳系与星际空间的分界。更远处、约 2,000 个天文单位（AU）以外，是太阳系最外层的假想区域——奥尔特云（Oort cloud）。它被认为是长周期彗星（long-period comets）的来源，其范围可能延伸至太阳系的引力界限（Hill sphere）边缘，约178,000–227,000 AU（2.81–3.59 光年），在此处太阳系的引力势与银河系的引力势达到平衡。目前，太阳系正穿越一个被称为本地星际云（Local Cloud）的星际介质区域。距离太阳系最近的恒星是比邻星（Proxima Centauri），距离约269,000 AU（4.25 光年）。太阳系与比邻星都位于本地泡（Local Bubble）内，这是银河系中一个直径约 1,000 光年（ly）的相对稀薄的区域。
\subsection{定义}
太阳系包括太阳以及所有受其引力约束并围绕其运行的天体。\(^\text{[15][16][17]}\)

国际天文学联合会（IAU）将太阳系描述为：由太阳引力所束缚的一切天体，包括太阳本身、八大行星，以及环绕太阳运行的其他天体。\(^\text{[11]}\)美国国家航空航天局（NASA）则将太阳系定义为一个行星系统（planetary system），其中包括太阳及所有环绕它运行的天体。\(^\text{[12]}\)

当“solar system”一词不作为专有名词、且首字母不大写时，它可以指太阳系本身，也可泛指任何类似太阳系的行星系统。\(^\text{[15]}\)
\subsection{形成与演化}
\subsubsection{过去}
\begin{figure}[ht]
\centering
\includegraphics[width=8cm]{./figures/0dc12bd0e5f767c3.png}
\caption{} \label{fig_Solar_2}
\end{figure}
太阳系形成于至少约46.68 亿年前，源自一片巨大分子云中某个区域的引力坍缩。\(^\text{[b]}\)这片最初的云可能横跨数光年，并极有可能同时孕育了多颗恒星。\(^\text{[19]}\)如同其他分子云一样，它的主要成分是氢气，其次是氦气，并含有少量由前几代恒星通过核聚变产生的较重元素。\(^\text{[20]}\)

当原太阳星云\(^\text{[20]}\)开始坍缩时，角动量守恒使其旋转速度不断加快。中心区域由于聚集了大部分质量，变得比周围更炽热、更致密。\(^\text{[16]}\)随着收缩的星云自转加速，它逐渐扁平化为一个直径约 200 天文单位（AU）的原行星盘\(^\text{[19][21]}\)，其中心形成了一个炽热而致密的原恒星。\(^\text{[22][23]}\)行星由盘中的物质通过吸积（accretion）过程形成，\(^\text{[24]}\)其中尘埃和气体相互引力吸引，逐渐聚合成更大的天体。在早期太阳系中，可能存在数百个原行星（protoplanets），但它们大多互相合并、被摧毁或被抛射出系统，最终仅留下行星、矮行星以及少量残余小天体。\(^\text{[25][26]}\)

在靠近太阳的内太阳系，即霜线（frost line）以内、甚至烟尘线（soot line）以内，\(^\text{[27]}\)由于高温使除金属与硅酸盐外的物质难以以固体形式存在，因此这一带形成的行星主要由岩石构成，即：水星、金星、地球与火星。由于这些耐高温的物质仅占太阳星云极小的比例，类地行星因此无法增长得太大。\(^\text{[25]}\)

四颗巨行星——木星、土星、天王星与海王星——在霜线之外形成，该区域位于火星与木星轨道之间，温度足够低，使挥发性冰类化合物可以保持固态。这些冰的丰度远高于类地行星形成所用的金属与硅酸盐，这使得巨行星能长得足够大，从而捕获大量的氢与氦 —— 这两种元素是宇宙中最轻、最丰富的。\(^\text{[25]}\)未能成为行星的残余碎片则聚集在小行星带、柯伊伯带与奥尔特云等区域。\(^\text{[25]}\)

约五千万年后，原恒星中心的氢气压与密度足够高，开始发生热核聚变反应。\(^\text{[28]}\)随着核心中氦的积累，太阳逐渐变得更亮；\(^\text{[29]}\)在主序星生命早期，其亮度仅为如今的约 70\%。\(^\text{[30]}\)温度、反应速率、压力与密度不断上升，直至达到静水平衡（hydrostatic equilibrium）：即热压与引力达到平衡。此时，太阳成为一颗主序星。\(^\text{[31]}\)太阳风从太阳表面喷射，形成了日球层（heliosphere），并将原行星盘中剩余的气体与尘埃吹散至星际空间。\(^\text{[29]}\)

在原行星盘消散后，尼斯模型（Nice model）提出：气体巨行星与微行星之间的引力相互作用导致行星迁移至不同轨道。这引发了整个系统的动力学不稳定，使微行星被散射，最终将气体巨行星带入如今的位置。在这一时期，大调车假说（Grand Tack Hypothesis）认为：木星的一次最终向内迁移扰乱了小行星带的大部分物质，从而引发了晚期重轰炸（Late Heavy Bombardment）。\(^\text{[32][33]}\)
\subsubsection{现在与未来}
太阳系目前处于一个相对稳定、缓慢演化的状态，所有天体都沿着受引力束缚的轨道绕太阳运行。\(^\text{[34]}\)尽管太阳系在数十亿年来一直较为稳定，但从技术角度上看，它是一个混沌系统，未来可能受到扰动。在未来数十亿年内，存在极小概率会有其他恒星穿越太阳系。这可能导致系统的轻微失稳，在数百万年后造成行星被逐出、相互碰撞或坠入太阳，但最有可能的结果仍是：太阳系将大体保持如今的样貌。\(^\text{[35]}\)
\begin{figure}[ht]
\centering
\includegraphics[width=8cm]{./figures/080701ee703e48af.png}
\caption{当前太阳与其在红巨星阶段达到峰值时的体积对比} \label{fig_Solar_3}
\end{figure}
太阳的主序阶段（main-sequence phase）——从开始到结束——将持续约 100 亿年，而此后所有阶段（从主序末期到白矮星前的各阶段）加起来的寿命仅约 20 亿年。\(^\text{[36]}\)太阳系将在此期间大体保持如今的状态，直到太阳核心中的氢完全转化为氦为止—— 这将发生在大约 50 亿年后。届时，太阳的主序生命将走到尽头。  

那时，太阳的核心将因引力进一步收缩，氢聚变将仅在包围惰性氦核的壳层（shell）中进行，其总能量输出将超过当前水平。太阳的外层将剧烈膨胀，直径将达到目前的约260 倍，并演化为一颗红巨星（Red Giant）。由于表面积大幅增加，其表面温度将下降，在最冷时约为2,600 K（4,220 °F），远低于主序阶段的温度。\(^\text{[36]}\)

膨胀后的太阳预计将汽化水星与金星，并使地球与火星变得不可居住（甚至可能完全摧毁地球）。\(^\text{[37][38]}\)最终，太阳核心温度将升高到足以点燃氦聚变反应；太阳将在核心燃烧氦，但这一阶段仅持续其氢燃烧阶段的极小一部分时间。由于太阳的质量不足以启动更重元素的聚变，核心中的核反应将逐渐衰竭并停止。太阳的外层将在这一过程中被抛射到太空中，留下一个致密的白矮星（White Dwarf）—— 其质量约为太阳原始质量的一半，但体积仅有地球大小。\(^\text{[36]}\)被抛射的外层气体可能形成一个行星状星云（Planetary Nebula），将太阳形成时的部分物质—— 此时已被碳等重元素所“富集”—— 重新送回星际介质（Interstellar Medium）。\(^\text{[39][40]}\)
\subsection{一般特征}
\begin{figure}[ht]
\centering
\includegraphics[width=8cm]{./figures/7b113d47a833549c.png}
\caption{这是一张经过色彩增强的从月球拍摄的太阳系多种成分的照片。左下角的三个光点自左向右依次为行星——土星、火星与水星；画面中央，太阳的日冕（corona）从月球漆黑的边缘上方升起，而月球右侧被地球反照光（earthshine）所照亮。} \label{fig_Solar_4}
\end{figure}
天文学家有时会将太阳系的结构划分为若干独立的区域。内太阳系（Inner Solar System）包括水星、金星、地球、火星以及位于小行星带中的天体；外太阳系（Outer Solar System）包括木星、土星、天王星、海王星以及位于柯伊伯带中的天体。\(^\text{[41]}\)自从发现柯伊伯带以来，太阳系的最外层部分被视为一个独立的区域，由位于海王星之外的天体组成。\(^\text{[42]}\)
\subsubsection{组成}
太阳系的主要成分是太阳，一颗G型主序星（G-type main-sequence star），其质量占整个系统已知总质量的99.86\%，在引力上完全主导着整个系统。\(^\text{[43]}\)围绕太阳运行的四个最大天体——四大行星——占据剩余质量的99\%，其中仅木星和土星就合计超过90\%。太阳系的其他所有天体（包括四颗类地行星、矮行星、卫星、小行星与彗星等）总共的质量还不到太阳系总质量的0.002\%。\(^\text{[h]}\)

太阳的成分大约由 98\%的氢与氦组成，\(^\text{[47]}\)木星与土星的成分也大体相同。\(^\text{[48][49]}\)太阳系中存在一个由早期太阳的热量与光压形成的成分梯度：距离太阳较近、受热与光压影响更强的天体，主要由熔点较高的元素构成；而距离太阳较远的天体，则主要由熔点较低的物质组成。\(^\text{[50]}\)在太阳系中，那条易挥发物质能够凝聚的界限称为霜线（frost line），其位置大约在距离太阳约地球距离的五倍处。\(^\text{[5]}\)
\subsubsection{轨道}
行星与其他围绕太阳运行的大型天体的轨道大多位于太阳系的不变平面（invariable plane）附近，地球的轨道（称为黄道面，ecliptic）也接近这一平面；其中木星的轨道与该平面的夹角最小，仅为 0.3219°。\(^\text{[51]}\)而较小的冰质天体（例如彗星）则常以更大的倾角绕此平面运行。\(^\text{[52][53]}\)太阳系中的大多数行星都拥有自己的“次级系统”，即被称为卫星（moons）的天然卫星所环绕。所有最大型的天然卫星都呈同步自转（synchronous rotation），始终以一面对着其母行星。四大行星（木星、土星、天王星、海王星）皆拥有行星环（planetary rings）—— 由无数微小颗粒组成的薄环盘，这些颗粒共同绕行行星。\(^\text{[54]}\)

由于太阳系形成的原因，行星与绝大多数其他天体的公转方向与太阳自转方向一致，即从地球北极上方观察为逆时针方向（counter-clockwise）。\(^\text{[55]}\)不过也有例外，例如哈雷彗星（Halley’s Comet）。\(^\text{[56]}\)多数较大的卫星以顺行方向（prograde direction）绕行行星，即与母行星的自转方向一致；而海王星的卫星海卫一（Triton）是已知最大的一颗逆行（retrograde）运行的卫星。\(^\text{[57]}\)大多数大型天体的自转方向相对于其公转方向也是顺行的，但金星（Venus）的自转方向为逆行。\(^\text{[58]}\)

在良好的近似下，开普勒行星运动定律（Kepler’s laws of planetary motion）可以描述天体绕太阳的轨道运动。\(^\text{[59]: 433–437}\) 这些定律指出：每个天体沿着一个椭圆轨道运行，而太阳位于该椭圆的一个焦点上，这导致天体在公转过程中与太阳的距离不断变化。天体距太阳最近的位置称为近日点（perihelion），最远的位置称为远日点（aphelion）。\(^\text{[60]: 9–6 }\)除水星外，行星的轨道几乎为圆形，而许多彗星、小行星和柯伊伯带天体则具有高度椭圆的轨道。开普勒定律只考虑太阳对轨道天体的引力作用，未考虑天体之间的相互引力。在人类的时间尺度上，这些扰动可通过数值模型加以修正，\(^\text{[60]: 9–6}\) 但在数十亿年的时间尺度上，整个行星系统可能会发生混沌变化。\(^\text{[61]}\)

太阳系的角动量（angular momentum）衡量了系统内所有运动成分的总轨道与自转动量。\(^\text{[62]}\)虽然太阳的质量占据绝对主导，但其仅贡献了大约2\% 的角动量。\(^\text{[63][64]}\)行星——主要是木星——由于其质量、轨道与距离太阳的组合，占据了系统绝大部分的角动量，而彗星可能也提供了显著的次要贡献。\(^\text{[63]}\)
\subsubsection{距离与尺度}
\begin{figure}[ht]
\centering
\includegraphics[width=8cm]{./figures/5dbc152db58202da.png}
\caption{太阳系中各行星相对轨道距离的可视化示意图——以压缩矩形形式呈现。} \label{fig_Solar_5}
\end{figure}
太阳的半径为0.0047 AU（70万千米；40万英里）。\(^\text{[65]}\)因此，太阳的体积仅占一个半径等于地球轨道半径的球体体积的0.00001\%（即 1/10^7），而地球的体积大约仅为太阳的百万分之一（10^{-6}）。最大行星木星距离太阳5.2 AU，半径为71,000 千米（0.00047 AU；44,000 英里）；最远的行星海王星则距离太阳30 AU。\(^\text{[49][66]}\)

除少数例外外，行星或行星带距离太阳越远，其轨道间距也越大。例如：金星比水星远约 0.33 AU；而土星比木星远4.3 AU；海王星则比天王星远10.5 AU。天文学家曾尝试寻找这些轨道间距的规律，如提丢斯–波得定律（Titius–Bode law）\(^\text{[67]}\)以及开普勒（Johannes Kepler）基于柏拉图立体（Platonic solids）的模型\(^\text{[68]}\)，但随着新的天体被发现，这些假说已被证伪。\(^\text{[69]}\)

太阳系比例模型（Solar System Scale Models）一些太阳系模型尝试以人类直观方式表现其中的相对尺度。 有的模型较小（甚至可为机械式，称为行星仪 orrery），也有的规模宏大，跨越整个城市或地区。\(^\text{[70]}\)目前已知最大的比例模型是瑞典的“瑞典太阳系模型（Sweden Solar System）”，它以斯德哥尔摩的阿维奇竞技场（Avicii Arena）—— 一座直径 110 米（361 英尺）的穹顶建筑——代表太阳。按照比例：木星是一颗直径 7.5 米（25 英尺）的球体，位于斯德哥尔摩阿兰达机场（40 千米 / 25 英里 之外）；而最远的天体 —— 赛德娜（Sedna）—— 则是一颗直径仅 10 厘米（4 英寸）的球体，位于吕勒奥（Luleå），距“太阳” 912 千米（567 英里）。\(^\text{[71][72]}\)在这一比例下，最近恒星比邻星（Proxima Centauri）的距离大约是地月距离的八倍。

若将太阳至海王星的距离缩放为100 米（330 英尺），则太阳的直径约为 3 厘米（1.2 英寸）—— 约为高尔夫球直径的三分之二；四大行星的直径均不超过 3 毫米（0.12 英寸）；而地球及其他类地行星的直径在该比例下将小于跳蚤（约 0.3 毫米 / 0.012 英寸）。\(^\text{[73]}\)
\begin{figure}[ht]
\centering
\includegraphics[width=8cm]{./figures/57f34943e92d54c3.png}
\caption{行星之间距离的比较图，白色横条表示轨道变化范围。行星的大小未按比例绘制。} \label{fig_Solar_6}
\end{figure}
\subsubsection{宜居性}
\begin{figure}[ht]
\centering
\includegraphics[width=8cm]{./figures/cc03a791c33b51ee.png}
\caption{太阳系与 TRAPPIST-1 恒星系宜居带的比较图。TRAPPIST-1 是一颗超冷红矮星，已知其周围有七颗类地行星以稳定轨道环绕运行。} \label{fig_Solar_7}
\end{figure}
\begin{figure}[ht]
\centering
\includegraphics[width=8cm]{./figures/caa7fcdf1175a2ac.png}
\caption{不同恒星温度下宜居带的比较图，示例包括若干已知系外行星以及地球、火星和金星。从上到下依次为：F 型主序星、黄矮星（G 型主序星）、橙矮星（K 型主序星）、典型红矮星，以及超冷矮星。} \label{fig_Solar_8}
\end{figure}
太阳系的宜居带（habitability zone）通常被认为位于内太阳系靠近地球的区域，在那里由于太阳辐射的作用，大气层中的液态水可以稳定存在。\(^\text{[74]}\)

除了太阳能外，太阳系中能够支持生命存在的主要特征是日球层（heliosphere）以及各行星自身的磁场（对于那些拥有磁场的行星而言）。这些磁场在一定程度上屏蔽了来自星际空间的高能粒子——宇宙射线。星际介质中宇宙射线的密度以及太阳磁场的强度会在极长的时间尺度上变化，因此太阳系中宇宙射线的渗透水平也随之变化，但其变化幅度目前仍未知。\(^\text{[75]}\)

太阳系中的宜居性不仅仅取决于表面条件或太阳辐射环境，因为在若干太阳系天体中可能存在的地下海洋中也可能具备生命存在的潜力，\(^\text{[76]}\)甚至在某些行星（尤其是金星）的云层中，也可能存在可居住层。\(^\text{[77]}\)
\subsubsection{与太阳系外行星系统的比较}
对开普勒（Kepler）望远镜观测数据的分析表明，银河系中被观测到的行星系统大致可分为三类：“相似型（similar）”系统——行星大小相近、间距相似且轨道高度圆形；“有序型（ordered）”系统——行星的质量随着与恒星距离的增加而增大；“混合型（mixed）”系统——行星质量分布无明显规律。太阳系属于“有序型”系统，占已观测系统的约37\%。然而，“相似型”系统为多数，占59\%，而“混合型”系统仅占4\%。\(^\text{[78]}\)

与许多太阳系外行星系统相比，太阳系的一个显著特征是 —— 缺乏位于水星轨道以内的行星。\(^\text{[79][80]}\)太阳系中也没有“超级地球”（质量为地球的1至10倍的行星），\(^\text{[79]}\)尽管假设中的“第九行星”（Planet Nine）若存在，可能是一颗位于太阳系边缘的超级地球。\(^\text{[81]}\)

另一个不寻常之处在于太阳系仅包含较小的类地行星和巨大的气体行星，而在其他行星系统中，中等大小的行星更为常见——包括岩质型与气态型——因此在地球与海王星（半径约为地球的3.8倍）之间的“尺寸空档”在太阳系中非常明显。由于许多超级地球的轨道都比水星更靠近其母星，有一种假说认为所有行星系统最初都有许多靠近恒星的行星，而随后它们通过一系列碰撞事件合并成更少但更大的行星。然而在太阳系中，这些碰撞反而导致行星的破坏与被抛出系统。\(^\text{[79][82]}\)

太阳系行星的轨道几乎为圆形。与许多其他系统相比，它们的轨道偏心率要小得多。\(^\text{[79]}\)尽管有尝试将这种现象部分归因于径向速度法的观测偏差，以及长期多行星相互作用的动力学效应，但其确切原因仍未确定。\(^\text{[79][83]}\)
\subsubsection{太阳}
\begin{figure}[ht]
\centering
\includegraphics[width=8cm]{./figures/cc5da4e20b80c9b8.png}
\caption{真实白光下的太阳（The Sun in true white color）} \label{fig_Solar_9}
\end{figure}
太阳（The Sun）是太阳系的中心恒星，也是其最为巨大的组成部分。它拥有巨大的质量（约为地球质量的 332,900 倍），\(^\text{[84]}\)占太阳系总质量的 99.86\%，\(^\text{[85]}\)这一庞大的引力使其核心达到足以维持氢聚变为氦的高温高压。\(^\text{[86]}\)这一聚变过程释放出极其巨大的能量，主要以电磁辐射的形式向外辐射，其辐射峰值位于可见光波段。\(^\text{[87][88]}\)

由于太阳在核心中进行氢聚变，因此它是一颗主序星。更具体地说，它属于G2 型主序星（G2V star），其中类型编号指其有效温度。主序星温度越高，其光度越强，但寿命越短。太阳的温度介于最热恒星与最冷恒星之间。比太阳更明亮、更炽热的恒星相对罕见；而显著更暗、更冷的恒星——即所谓的红矮星——约占银河系中所有核聚变恒星的 75\%。\(^\text{[89]}\)

太阳属于第一族恒星（Population I star），形成于银河系的旋臂中。与银河系核球与晕中的老年第二族恒星相比，它含有更多比氢与氦更重的元素（在天文学中称为“金属”）。\(^\text{[90]}\)这些较重元素是在更早一代恒星的核心与爆炸中形成的，因此宇宙必须先经历第一代恒星的死亡，才能被这些重元素所富集。最古老的恒星几乎不含金属，而后期形成的恒星则金属丰度更高。人们认为这种较高的金属丰度对太阳形成行星系统至关重要，因为行星正是由这些“金属”物质的吸积所形成的。\(^\text{[91]}\)

由太阳磁层主导的空间区域称为日球层（heliosphere），它覆盖了太阳系的大部分区域。除了光之外，太阳还不断释放带电粒子流（等离子体），称为太阳风（solar wind）。太阳风以每小时 90 万公里（56 万英里）至 288 万公里（179 万英里）的速度向外扩散，\(^\text{[92]}\)充满了太阳系天体之间的真空空间，形成了一种稀薄、含尘的介质，称为行星际介质（interplanetary medium），其范围至少延伸至 100 天文单位。\(^\text{[93]}\)

太阳表面的活动，如太阳耀斑（solar flares）与日冕物质抛射（coronal mass ejections, CME），会扰动整个日球层，产生 空间天气（space weather）并引发 地磁风暴（geomagnetic storms）。\(^\text{[94]}\)日冕物质抛射及类似事件会将磁场与大量物质从太阳表面喷射出来。这些磁场与物质与地球磁场相互作用，将带电粒子引导至地球高层大气，其相互作用形成极光（aurora），可在地磁极附近观测到。\(^\text{[95]}\)日球层中最大的稳定结构是 日球电流片（heliospheric current sheet），它是由太阳自转磁场作用于行星际介质形成的螺旋形结构。\(^\text{[96][97]}\)
\subsection{内太阳系}
内太阳系是指包含类地行星与小行星的区域。\(^\text{[98]}\)该区域内的天体主要由硅酸盐和金属组成，\(^\text{[99]}\)并且它们距离太阳相对较近；整个区域的半径小于木星与土星轨道之间的距离。此区域位于冰霜线（frost line）之内，该冰霜线距离太阳略小于 5 个天文单位（AU）。\(^\text{[52]}\)
\subsubsection{类地行星}
\begin{figure}[ht]
\centering
\includegraphics[width=8cm]{./figures/d502b9668b686c61.png}
\caption{四颗类地行星：水星、金星、地球与火星（The four terrestrial planets Mercury, Venus, Earth and Mars）} \label{fig_Solar_10}
\end{figure}
四颗类地行星（或称内行星）具有致密的岩质组成，卫星稀少或完全没有，也不存在行星环系统。它们主要由耐高温的矿物（如硅酸盐）组成，这些矿物形成了行星的地壳与地幔；而其核心则主要由铁与镍等金属构成。四颗行星中的三颗——金星、地球和火星——拥有足够厚的气体包层，能够形成天气系统；所有类地行星都具有陨石坑和构造地貌，如裂谷与火山。\(^\text{[100]}\)
\begin{itemize}
\item 水星（距离太阳 0.31–0.59 AU）\(^\text{[D 6]}\)是太阳系中最小的行星。其表面呈灰色，拥有由逆冲断层形成的广阔断崖（rupes）系统，以及由撞击事件残迹形成的明亮射纹系统。\(^\text{[101]}\)表面温度变化剧烈——赤道区域夜间温度可达 −170 °C（−270 °F），而白天则高达 420 °C（790 °F）。在过去，水星曾有火山活动，形成了类似月球的光滑玄武质平原。\(^\text{[102]}\)水星可能拥有硅酸盐地壳与巨大的铁核。\(^\text{[103][104]}\)它的大气极其稀薄，仅由太阳风粒子与喷射出的原子组成。\(^\text{[105]}\)水星没有天然卫星。\(^\text{[106]}\)
\item 金星（距离太阳 0.72–0.73 AU）\(^\text{[D 6]}\)拥有高度反射的白色大气，主要由二氧化碳组成。其表面气压约为地球海平面气压的 90 倍。\(^\text{[107]}\)金星表面温度超过 400 °C（752 °F），主要原因是其大气中温室气体含量极高。\(^\text{[108]}\)该行星缺乏可抵御太阳风剥蚀的磁场，这表明其大气可能由火山活动持续补充。\(^\text{[109]}\)其表面显示出广泛的火山活动证据，属于“停滞盖板构造”（stagnant-lid tectonics）。\(^\text{[110]}\)金星没有天然卫星。\(^\text{[106]}\)
\item 地球（距离太阳 0.98–1.02 AU）[D 6]是宇宙中目前已知唯一存在生命与地表液态水的天体。\(^\text{[111]}\)地球大气中含有 78\% 的氮气与 21\% 的氧气，这正是生命活动的结果。\(^\text{[112][113]}\)地球拥有复杂的气候与天气系统，不同气候带之间差异显著。\(^\text{[114]}\)地表主要由绿色植被、荒漠与白色冰盖组成。\(^\text{[115][116][117]}\)地球的板块构造塑造了大陆形态。\(^\text{[102]}\)其行星磁层（magnetosphere）能屏蔽辐射，减少大气流失，从而维持生命宜居性。\(^\text{[118]}\)
\item 月球是地球唯一的天然卫星。\(^\text{[119]}\)其直径约为地球的四分之一。\(^\text{[120]}\)表面覆盖着极细的风化层（regolith），并布满撞击坑。\(^\text{[121][122]}\)月球上的大暗斑——“月海”（maria）——由古代火山活动形成。\(^\text{[123]}\)其大气极其稀薄，仅为部分真空，粒子密度低于每立方厘米 10^7。\(^\text{[124]}\)
\item 火星（距离太阳 1.38–1.67 AU）[D 6]其半径约为地球的一半。\(^\text{[125]}\)由于火星土壤中的氧化铁，其表面呈红色。\(^\text{[126]}\)极地地区覆盖着由水和二氧化碳组成的白色冰盖。\(^\text{[127]}\)火星大气主要由二氧化碳组成，表面气压仅为地球的 0.6\%，但仍足以产生一些天气现象。\(^\text{[128]}\)火星年为 687 个地球日，其表面温度变化剧烈，范围从 −78.5 °C（−109.3 °F）至 5.7 °C（42.3 °F）。表面布满火山与裂谷，并富含多种矿物。\(^\text{[129][130]}\)火星内部结构高度分化，并在约 40 亿年前失去了磁层。\(^\text{[131][132]}\)火星拥有两颗微小卫星：\(^\text{[133]}\)
\item 火卫一（Phobos）是火星的内卫星。其形状不规则，平均半径为 11 km（7 mi）。表面反射率极低，布满撞击坑。\(^\text{[D 7][134]}\)其中最大的撞击坑为 Stickney，半径约 4.5 km（2.8 mi）。\(^\text{[135]}\)
\item 火卫二（Deimos）是火星的外卫星。与火卫一相似，形状不规则，平均半径 6 km（4 mi），表面反光度同样极低。[D 8][D 9]然而，火卫二的表面比火卫一更为光滑，因为风化层部分覆盖了撞击坑。\(^\text{[136]}\)
\end{itemize}
\subsubsection{小行星}
\begin{figure}[ht]
\centering
\includegraphics[width=8cm]{./figures/e8184c45bcc8665a.png}
\caption{内太阳系概览 —— 直至木星轨道（Overview of the inner Solar System up to Jupiter's orbit）} \label{fig_Solar_11}
\end{figure}
除最大的小行星谷神星（Ceres）外，其他小行星均被归类为“太阳系小天体”（small Solar System bodies）。它们主要由碳质物质、耐高温的岩石与金属矿物组成，部分含有冰。\(^\text{[137][138]}\)其尺寸范围从数米至数百公里不等。许多小行星依据其轨道特征被划分为不同的小行星族群（groups）与家族（families）。部分小行星拥有绕其运行的天然卫星，即所谓“次级小行星”——绕更大行星体运行的小行星。\(^\text{[139]}\)
\begin{itemize}
\item 穿越水星轨道的小行星（Mercury-crossing asteroids）指近日点轨道位于水星轨道以内的小行星。截至目前已知至少有 362 颗，它们包括太阳系中距离太阳最近的天体。\(^\text{[140]}\)目前尚未发现任何位于水星轨道与太阳之间的“火神星族”（vulcanoids）。\(^\text{[141][142]}\)截至 2024 年，仅发现一颗完全在金星轨道以内运行的小行星——594913ꞌAylóꞌchaxnim。\(^\text{[143]}\)
\item 穿越金星轨道的小行星（Venus-crossing asteroids）指轨道与金星轨道相交的小行星。截至 2015 年，共记录 2,809 颗。\(^\text{[144]}\)
\item 近地小行星（Near-Earth asteroids）其轨道相对接近地球轨道。\(^\text{[145]}\)部分被列为潜在危险天体（potentially hazardous objects），因为它们未来有可能与地球发生碰撞。\(^\text{[146][147]}\)截至 2024 年，已知的近地小行星数量超过 37,000 颗。\(^\text{[148]}\)一些绕太阳运行的流星体（meteoroids）在撞击地球前足够大，能够被直接观测与追踪。如今普遍认为，过去的撞击事件在塑造地球的地质与生物演化历史中起到了重要作用。\(^\text{[149]}\)
\item 穿越火星轨道的小行星（Mars-crossing asteroids）指近日点大于 1.3 AU、轨道穿越火星轨道的小行星。\(^\text{[150]}\)截至 2024 年，美国 NASA 已确认 26,182 颗火星交叉小行星。\(^\text{[144]}\)
\end{itemize}
\textbf{小行星带}

小行星带位于距离太阳约 2.3 至 3.3 个天文单位（AU）的环状区域，其位置介于火星与木星轨道之间。一般认为，它由太阳系形成时期未能聚合成行星的残余物质组成，这种失败主要是由于木星强大的引力扰动所致。\(^\text{[151]}\)小行星带中含有数以万计、甚至可能数百万计直径超过一公里的天体。\(^\text{[152]}\)尽管数量庞大，但其总体质量不太可能超过地球质量的千分之一。\(^\text{[46]}\)小行星带的物质分布极为稀疏，探测器在穿越该区域时通常不会遭遇任何碰撞或损伤。\(^\text{[153]}\)
\begin{figure}[ht]
\centering
\includegraphics[width=8cm]{./figures/7ea9a2828e334e64.png}
\caption{四大小行星：谷神星（Ceres）、灶神星（Vesta）、智神星（Pallas）与健康神星（Hygiea）。目前仅有谷神星与灶神星被航天器造访过，因此只有它们拥有详细的表面影像。} \label{fig_Solar_12}
\end{figure}
以下是小行星带中三大天体的简介。它们都被认为是相对完整的原行星（protoplanets），即行星在完全形成前的早期阶段（参见 特殊小行星列表）：\(^\text{[154][155][156]}\) 
\begin{itemize}
\item 谷神星（Ceres，2.55–2.98 AU）谷神星是小行星带中唯一的矮行星。\(^\text{[157]}\)它是该区域中最大的天体，直径约为 940 公里（580 英里）。\(^\text{[158]}\)其表面由碳质物质、冻结水与含水矿物的混合物组成。\(^\text{[159][160]}\)表面存在过去低温火山活动（cryovolcanism）的迹象，这类活动会喷发出水等挥发性物质，在表面形成亮斑区域。\(^\text{[161]}\)谷神星拥有极其稀薄的水汽大气层，但实际上几乎与真空无异。\(^\text{[162]}\) 
\item 灶神星（Vesta，2.13–3.41 AU）灶神星是小行星带中第二大天体。\(^\text{[163]}\)其碎片仍以“灶神星族”（Vesta family）的形式存在，\(^\text{[164]}\)并且在地球上发现了大量来源于它的 HED 陨石。\(^\text{[165]}\)灶神星表面以玄武质与变质物质为主，其密度高于谷神星。\(^\text{[166]}\)表面具有两个巨大的撞击坑：Rheasilvia 与 Veneneia。\(^\text{[167]}\) 
\item 智神星（Pallas，2.15–2.57 AU）智神星是小行星带中第三大天体。\(^\text{[163]}\)它拥有自己的“智神星族”。\(^\text{[164]}\)由于尚未有航天器造访智神星，人们对其了解有限，\(^\text{[168]}\)但推测其表面主要由硅酸盐物质构成。\(^\text{[169]}\)
\end{itemize}
希尔达族小行星（Hilda asteroids）这类小行星与木星呈 3:2 共振关系，即它们每绕太阳三周，木星绕行两周。\(^\text{[170]}\)它们分布在木星轨道与主小行星带之间的三个互相关联的聚集区中。  

特洛伊小行星（Trojans）这类天体位于其他行星的引力稳定拉格朗日点处：\(L_4\)点位于行星轨道前方 60°，\(L_5\)点位于其轨道后方 60°。\(^\text{[171]}\)除水星外，所有行星都已知至少拥有一个特洛伊小行星。\(^\text{[172][173][174]}\)木星的特洛伊小行星数量大约与整个小行星带相当。\(^\text{[175]}\)在木星之后，海王星拥有第二多的特洛伊天体，目前确认的数量为 28 个。\(^\text{[176]}\)
\subsection{外行星}
\begin{figure}[ht]
\centering
\includegraphics[width=6cm]{./figures/520245b686f396f2.png}
\caption{外行星——木星、土星、天王星与海王星，与右下角的内行星——地球、金星、火星和水星——的对比图。} \label{fig_Solar_13}
\end{figure}
四颗外行星，又称巨行星（giant planets）或类木行星（Jovian planets），共同构成了围绕太阳运行天体质量的 99\%。\(^\text{[h]}\)所有四颗巨行星都有多颗卫星与行星环系统，但从地球上肉眼仅能轻易观测到土星的环。\(^\text{[100]}\)木星与土星主要由熔点极低的气体（如氢、氦与氖）组成，\(^\text{[177]}\)因此被称为气体巨行星（gas giants）。\(^\text{[178]}\)天王星与海王星则为冰巨行星（ice giants），\(^\text{[179]}\)即在天文学意义上主要由“冰”组成 —— 包括熔点在数百开尔文以内的化合物（如水、甲烷、氨、硫化氢与二氧化碳）。\(^\text{[177][180]}\)冰质物质也是巨行星卫星与海王星轨道以外小天体的主要成分。\(^\text{[180][181]}\)  
\begin{itemize}
\item 木星（Jupiter, 4.95–5.46 AU）\(^\text{[D 6]}\)木星是太阳系中体积最大、质量最高的行星。其表面可见橙褐与白色云带，按照大气环流原理高速运动，并伴有如“大红斑（Great Red Spot）”及白色“卵形风暴”等巨大气旋。木星的磁层极为强大，足以偏转电离辐射并在两极产生极光。\(^\text{[182]}\)截至 2025 年，木星已确认有 97 颗卫星，大致可分为三类：  
\item 阿玛尔忒亚群（Amalthea group）：包含 Metis、Adrastea、Amalthea 与 Thebe，它们的轨道显著近于木星，且其物质构成了木星微弱的行星环。\(^\text{[183][184]}\)  
\item 伽利略卫星（Galilean moons）：包含 Ganymede（盖尼米得）、Callisto（卡利斯托）、Io（木卫一）与 Europa（欧罗巴），它们是木星最大、最具行星特征的卫星。\(^\text{[185]}\) 
\item 不规则卫星（Irregular satellites）：体积较小，轨道更远且偏心率更高。\(^\text{[186]}\) 
\end{itemize}
\begin{itemize}
\item 土星（Saturn, 9.08–10.12 AU）\(^\text{[D 6]}\)土星最显著的特征是其围绕赤道的壮丽环系统，由冰与岩石微粒组成。与木星类似，它主要由氢与氦构成。\(^\text{[187]}\)土星南北两极存在巨大的六边形风暴，其直径超过地球。土星的磁层能产生微弱极光。截至 2025 年，土星已确认有 274 颗卫星，主要分为：  
\item 环内小卫星与牧羊卫星（Ring moonlets & shepherds）：它们分布在环内或其附近。小卫星仅能部分清理轨道尘埃，\(^\text{[188]}\)而牧羊卫星能完全清空轨道物质，形成可见间隙。\(^\text{[189]}\)  
\item 内层大卫星：包括 Mimas、Enceladus、Tethys 与 Dione。它们位于 E 环内，由水冰为主组成，内部结构已分化。\(^\text{[190]}\)  
\item 特洛伊卫星（Trojan moons）：Calypso 与 Telesto（Tethys 的特洛伊卫星），Helene 与 Polydeuces（Dione 的特洛伊卫星），这些小卫星与各自母卫星共享轨道。\(^\text{[191][192]}\)  
\item 外层大卫星：Rhea、Titan、Hyperion 与 Iapetus。\(^\text{[190]}\)其中 土卫六 Titan是太阳系中唯一拥有浓厚大气层的卫星。\(^\text{[193]}\) 
\item 不规则卫星：轨道更远、形状不规则，其中最大者为 Phoebe。\(^\text{[194]}\) 
\end{itemize}
\begin{itemize}
\item 天王星（Uranus, 18.3–20.1 AU）\(^\text{[D 6]}\)天王星独特地“横躺”在轨道上，其自转轴倾角超过 90°。这导致它经历极端的季节变化，两极交替指向与背离太阳。\(^\text{[195]}\)天王星外层呈黯淡的青绿色，但其内部气候仍存诸多谜团，如极低的内部热量与不规则的云层活动。  
截至 2025 年，天王星已确认 28 颗卫星，可分为三类：  
\item 内卫星：位于环系统内部，\(^\text{[196]}\)它们之间的距离极近，轨道呈混沌状态。\(^\text{[197]}\)  
\item 大卫星：Titania、Oberon、Umbriel、Ariel 与 Miranda。\(^\text{[198]}\)大多数由近乎等量的岩石与冰组成，唯独 Miranda 主要由冰构成。\(^\text{[199]}\)  
\item 不规则卫星：轨道更远且离心率更大。\(^\text{[200]}\) 
\end{itemize}
\begin{itemize}
\item 海王星（Neptune, 29.9–30.5 AU）\(^\text{[D 6]}\)海王星是目前已知距离太阳最远的行星。其外层大气呈暗青色，偶尔出现表面暗斑风暴。如同天王星，其大气现象仍未完全解释，包括异常高温的热层及磁层倾角高达 47°。截至 2025 年，海王星已确认 16 颗卫星，主要分为：  
\item 规则卫星（Regular satellites）： 轨道近圆形，接近赤道面。\(^\text{[194]}\)  
\item 不规则卫星（Irregular satellites）：轨道不规则。其一 —— 海卫一 Triton，是海王星最大卫星。它地质活动活跃，存在喷发氮气的间歇泉，并拥有稀薄的氮云大气层。\(^\text{[201][193]}\)
\end{itemize}
\subsubsection{半人马小行星}
半人马类天体是冰质的、类彗星的小天体，其轨道半长轴大于木星而小于海王星（约 5.5～30 个天文单位之间）。它们是曾经属于柯伊伯带（Kuiper Belt）或散射盘（Scattered Disc Objects, SDO）的天体，后被外行星的引力扰动而迁入更靠近太阳的区域，预计最终会演化为彗星或被逐出太阳系。\(^\text{[45]}\)大多数半人马类天体呈不活跃的类小行星状态，但部分天体会表现出彗星活动。例如首个被发现的半人马天体 —— 2060 Chiron（凯龙星），因在接近太阳时会形成彗发（coma），而被归类为彗星（95P）。\(^\text{[202]}\)已知最大的半人马天体是10199 Chariklo（查里克洛），直径约为 250 千米（160 英里），它是少数被发现拥有环系统的小行星之一。\(^\text{[203][204]}\)  
\subsection{海王星外区域}
在海王星轨道之外，是所谓的“海王星外区域”，其中包括环状的柯伊伯带（Kuiper Belt）—— 冥王星（Pluto）及其他多颗矮行星的所在地，以及与其重叠、相对于太阳系平面略倾斜并向更远处延伸的散射盘（Scattered Disc）。这一整片区域仍然几乎未被充分探测。它似乎主要由成千上万颗小型天体组成，其中最大者的直径仅为地球的五分之一，质量远小于月球。这些天体主要由岩石与冰构成。该区域有时被称为“太阳系的第三层区域（the third zone of the Solar System）”，环绕并延伸于内、外太阳系之外。\(^\text{[205]}\)
\subsubsection{柯伊伯带}
\begin{figure}[ht]
\centering
\includegraphics[width=6cm]{./figures/d56c747a4d526d4a.png}
\caption{柯伊伯带及其他小行星族群周围天体的分布图。图中 J、S、U、N 分别代表木星（Jupiter）、土星（Saturn）、天王星（Uranus）与海王星（Neptune）。} \label{fig_Solar_14}
\end{figure}
\begin{figure}[ht]
\centering
\includegraphics[width=6cm]{./figures/5f22641cc65a0277.png}
\caption{柯伊伯带天体的轨道分类图。  其中部分受轨道共振影响的天体族群已在图中标出。} \label{fig_Solar_15}
\end{figure}
柯伊伯带（Kuiper belt）是一条巨大的碎片环，类似于小行星带，但主要由冰质天体组成。\(^\text{[206]}\)它分布在距离太阳约 30 至 50 天文单位（AU）之间。该区域主要由小型太阳系天体构成，尽管其中最大的一些天体可能已经足够大，可以被归类为矮行星。\(^\text{[207]}\)估计柯伊伯带中直径超过 50 公里（30 英里）的天体数量超过十万个，但其总质量被认为仅为地球质量的十分之一，甚至百分之一。\(^\text{[45]}\)许多柯伊伯带天体拥有卫星，\(^\text{[208]}\)且大多数轨道相对于黄道平面有较大倾角（约 10°）。\(^\text{[209]}\)

柯伊伯带可大致分为“经典带”（classical belt）与“共振区”（resonant trans-Neptunian objects）。\(^\text{[206]}\)后者的轨道周期与海王星呈简单整数比，例如绕太阳运行两圈时海王星运行三圈，或一圈对应两圈。经典带由与海王星无共振的天体组成，范围大致从 39.4 AU 延伸至 47.7 AU。\(^\text{[210]}\)经典柯伊伯带的成员有时被称为“cubewanos”，名称来源于首个被发现的同类天体——1992 QB₁（现名 Albion），这些天体依然保持着接近原始的低偏心率轨道。\(^\text{[211]}\)

目前天文学界普遍认为，柯伊伯带中至少有五个天体是矮行星。\(^\text{[207][212]}\)同时，还有更多矮行星候选体正在等待观测确认。\(^\text{[213]}\)
\begin{itemize}
\item 冥王星（Pluto，29.7–49.3 AU）是柯伊伯带中已知最大的天体。其轨道偏心率较高，并相对于黄道面倾斜 17°。冥王星与海王星存在 2:3 的轨道共振关系，即冥王星绕太阳两圈时，海王星完成三圈。与冥王星轨道共振的柯伊伯带天体被称为“冥族天体”（plutinos）。\(^\text{[214]}\)冥王星拥有五颗卫星：卡戎（Charon）、冥卫一 Styx、冥卫二 Nix、冥卫三 Kerberos、冥卫四 Hydra。\(^\text{[215]}\)其中卡戎是最大的一颗，有时被认为与冥王星构成“矮行星双星系统”，因为两者的质心位于两体之外（即它们似乎彼此“互相绕转”）。
\item 奥库斯（Orcus，30.3–48.1 AU）与冥王星一样处于与海王星的 2:3 共振轨道，是继冥王星之后该共振族中最大的成员。\(^\text{[216]}\)其轨道偏心率与倾角与冥王星相似，但近日点与冥王星相差约 120°。因此，当冥王星位于近日点（最近在 1989 年）时，奥库斯位于远日点（最近在 2019 年），反之亦然。\(^\text{[217]}\)基于这种轨道“相反”关系，它被称为“反冥王星”（anti-Pluto）。\(^\text{[218][219]}\)奥库斯拥有一颗已知卫星——Vanth。\(^\text{[220]}\)
\item 哈梅亚（Haumea，34.6–51.6 AU）于 2005 年发现，\(^\text{[221]}\)目前处于与海王星 7:12 的暂时性轨道共振。[216]哈梅亚拥有一个环系统、两颗卫星（Hiʻiaka 与 Namaka），并且自转速度极快（约 3.9 小时一转），导致其形状被拉伸成椭球体。它属于一类具有相似轨道的碰撞族成员，说明数十亿年前，一次巨大的撞击从哈梅亚上剥离出大量碎片进入太空。\(^\text{[222]}\)
\item 鸟神星（Makemake，38.1–52.8 AU）虽小于冥王星，但却是“经典柯伊伯带”中已知最大的天体（即与海王星无共振关系的柯伊伯带天体）。鸟神星是柯伊伯带中仅次于冥王星的最亮天体，2005 年发现，2009 年正式命名。\(^\text{[223]}\)其轨道倾角达 29°，远高于冥王星。\(^\text{[224]}\)鸟神星拥有一颗已知卫星：S/2015 (136472) 1。\(^\text{[225]}\)
\item 夸欧尔（Quaoar，41.9–45.5 AU）是“经典柯伊伯带”中第二大已知天体，仅次于鸟神星。与鸟神星和哈梅亚相比，夸欧尔的轨道偏心率与倾角都要小得多。\(^\text{[216]}\)它拥有一个环系统和一颗已知卫星——Weywot。\(^\text{[226]}\)
\end{itemize}
\subsubsection{散射盘}
\begin{figure}[ht]
\centering
\includegraphics[width=6cm]{./figures/939833290b7e5420.png}
\caption{散射盘天体的轨道偏心率与轨道倾角，与经典柯伊伯带天体及共振柯伊伯带天体的比较图。} \label{fig_Solar_16}
\end{figure}
散射盘（Scattered disc）与柯伊伯带部分重叠，但向外延伸至约 500 天文单位（AU），被认为是短周期彗星的主要来源。散射盘天体被认为是在海王星早期向外迁移的过程中，受到其引力扰动而被抛入不规则轨道的。\(^\text{[227]}\)大多数散射盘天体的近日点位于柯伊伯带范围内，但远日点却远超柯伊伯带（有些距离太阳超过 150 AU）。这些天体的轨道相对于黄道面的倾角可高达 46.8°。\(^\text{[227]}\)一些天文学家认为散射盘仅是柯伊伯带的一个延伸区域，并将散射盘天体称为“散射柯伊伯带天体”（scattered Kuiper belt objects）。\(^\text{[228]}\)另一些学者则将半人马型天体（Centaurs）归为“向内散射的柯伊伯带天体”，与散射盘中的“向外散射天体”对应。\(^\text{[229]}\) 

目前，天文学界普遍认为散射盘中至少有两颗天体是矮行星：
\begin{itemize}
\item 厄里斯（Eris，38.3–97.5 AU）是目前已知最大的散射盘天体，也是质量最大的矮行星。厄里斯的发现引发了关于“行星定义”的争论，因为它的质量比冥王星大约高出 25\%，\(^\text{[230]}\)而直径几乎相当。它拥有一颗已知卫星 —— 纷争女神（Dysnomia）。与冥王星类似，厄里斯的轨道偏心率很高，近日点为 38.2 AU（与冥王星相当），远日点为 97.6 AU，且其轨道相对于黄道面倾斜 44°。\(^\text{[231]}\)
\item 龚龚（Gonggong，33.8–101.2 AU）是另一颗矮行星，其轨道与厄里斯相似，但与海王星存在 3:10 的轨道共振关系。\(^\text{[D 10]}\)它拥有一颗已知卫星 —— 相柳。\(^\text{[232]}\)
\end{itemize}
\subsubsection{极端跨海王星天体}
\begin{figure}[ht]
\centering
\includegraphics[width=6cm]{./figures/4522d4fe938a8038.png}
\caption{赛德娜（Sedna）、2012 VP113、Leleākūhonua（粉色）及其他极远天体（红色、棕色与青色）的当前轨道，以及假设中的第九行星（Planet Nine）的预测轨道（深蓝色）。} \label{fig_Solar_17}
\end{figure}
太阳系中有一些天体的轨道极为遥远，  
因此它们受到已知巨行星（如木星、土星、天王星、海王星）引力影响的程度远小于其他小行星族群。这些天体被称为极端跨海王星天体（Extreme Trans-Neptunian Objects，简称 ETNOs）。\(^\text{[233]}\)一般而言，ETNOs 的轨道半长轴至少在 150–250 天文单位（AU）之间。\(^\text{[233][234]}\)例如，541132 Leleākūhonua 绕太阳一周需约 32,000 年，  
其距离太阳的范围在 65–2000 AU 之间。\(^\text{[D 11]}\)  

天文学家将这一族群分为三类：散射型 ETNOs（Scattered ETNOs）其近日点约在 38–45 AU，轨道偏心率极高（>0.85）。与普通散射盘天体类似，它们很可能是早期被海王星引力散射形成的，并且仍然会与巨行星发生引力相互作用。 分离型 ETNOs（Detached ETNOs）其近日点大约在 40–45 AU 到 50–60 AU 之间，受海王星影响较小，但仍处于相对接近的范围内。赛德娜族或内奥尔特云天体（Sednoids / Inner Oort Cloud Objects）其近日点超过 50–60 AU，距离海王星太远，几乎不再受其引力影响。\(^\text{[233]}\)  

目前，只有一颗极端跨海王星天体被确认为矮行星：  
\begin{itemize}
\item 赛德娜（Sedna，76.2–937 AU）是首个被发现的极端跨海王星天体。它是一颗体积较大、呈红色的天体，绕太阳公转一周约需 11,400 年。其发现者迈克·布朗（Mike Brown，2003 年）认为，赛德娜不可能属于散射盘或柯伊伯带，因为其近日点距离过远，不可能受到海王星迁移的影响。\(^\text{[235]}\)“赛德娜族”（Sednoids）这一名称正是以它命名的。\(^\text{[233]}\)  
\end{itemize}
在部分 ETNO 的轨道中，天文学家发现其近日点分布存在统计聚集现象：它们最接近太阳的点集中在同一空间区域，且这些天体的轨道倾角也非常相似。\(^\text{[236][237][238]}\)部分天文学家推测，这可能是海王星外一颗大型行星的引力所致 —— 这一假设天体被称为第九行星（Planet Nine）。\(^\text{[239]}\)但也有研究者认为，这种轨道聚集可能仅仅是观测偏差或随机巧合所致。\(^\text{[240]}\)  
\subsubsection{奥尔特云}
\begin{figure}[ht]
\centering
\includegraphics[width=6cm]{./figures/4ddcebc02c5ee000.png}
\caption{艺术家绘制的奥尔特云（Oort Cloud）示意图：该区域仍位于太阳系引力势能范围之内。插图展示了位于其内部更近处的柯伊伯带（Kuiper Belt）。（图中天体的尺寸为便于观察而被夸大。）} \label{fig_Solar_18}
\end{figure}
奥尔特云（Oort Cloud）是一种假设存在的球形壳层结构，可能包含多达一万亿个冰质天体。它被认为是所有长周期彗星的来源。\(^\text{[241][242]}\)这些天体最初由太阳系内部形成，后来因与外行星的引力相互作用而被抛射到遥远的轨道上。奥尔特云中的天体运动极其缓慢，但会受到一些偶发事件的扰动，例如：与其他天体的碰撞；掠过恒星的引力效应；以及来自银河系的潮汐力（即所谓的“星系潮”）。\(^\text{[241][242]}\)由于其距离极为遥远、光度极低，以目前的观测技术尚无法直接探测到奥尔特云。\(^\text{[243]}\)  

根据理论模型，奥尔特云可能从距离太阳约 2,000 AU 延伸至 200,000 AU。而较低的估计则认为其外缘不会超过 50,000 AU。\(^\text{[244]}\)其主要质量集中在 3,000 至 100,000 AU 之间的区域。\(^\text{[245]}\)已知最远的天体（如西彗星（Comet West））其远日点大约位于 70,000 AU 处。\(^\text{[246]}\)  
\subsection{引力不稳定族群}
\subsubsection{流星体、流星与尘埃}
\begin{figure}[ht]
\centering
\includegraphics[width=6cm]{./figures/b7607936f8bdca18.png}
\caption{行星、黄道光与流星雨（图像左上方）} \label{fig_Solar_19}
\end{figure}
直径小于一米的固体天体通常被称为流星体（meteoroids），更小的（粒状）则称为微流星体（micrometeoroids），但两者的确切划分标准多年来一直存在争议。\(^\text{[247]}\)截至2017 年，国际天文学联合会（IAU）规定：直径约在 30 微米至 1 米 之间的固体天体统一称为流星体（meteoroids），并弃用了“微流星体”这一分类；而更小的颗粒则被简单称为尘埃粒子（dust particles）。\(^\text{[248]}\)

部分流星体来源于彗星或小行星的解体，另有少数来自行星表面撞击所产生的喷射碎片。大多数流星体由硅酸盐及重金属（如镍、铁）组成。\(^\text{[249]}\)当彗星经过太阳系时，其物质因蒸发或破碎形成一条由流星体构成的轨迹。当这些流星体进入行星大气层时，与空气摩擦发光，产生天空中明亮的划线现象，即流星（meteor）。如果一群流星体以平行轨迹同时进入大气层，它们在天空中看似从某一点“辐射”出来，这种现象称为流星雨（meteor shower）。\(^\text{[250]}\) 

太阳系内区存在黄道尘埃云（zodiacal dust cloud），在无光污染的暗夜中可见为黄道光（zodiacal light）。该尘埃云可能由小行星带中因行星引力扰动引发的碰撞形成，也有研究提出其部分物质可能源自火星。\(^\text{[251]}\)太阳系外区也存在宇宙尘埃云（cosmic dust cloud），范围约在 10 AU 至 40 AU 之间，被认为主要由柯伊伯带天体间的碰撞所产生。\(^\text{[252][253]}\)  
\subsubsection{彗星}
\begin{figure}[ht]
\centering
\includegraphics[width=6cm]{./figures/047b9ea4e9e01eef.png}
\caption{海尔–波普彗星尾部稀薄的等离子体（蓝色）与尘埃（白色）， 分别受太阳辐射压力与太阳风的作用而形成其形态。} \label{fig_Solar_20}
\end{figure}
彗星（Comets）是太阳系中的小型天体，通常直径只有数公里，主要由挥发性冰质物质组成。它们的轨道具有极高的离心率：近日点一般位于类地行星轨道内，而远日点则常远远超出冥王星轨道。当彗星进入内太阳系时，靠近太阳的高温使其冰质表面升华并电离，从而形成彗发（coma）—— 一条由气体与尘埃组成的长尾，常可肉眼观测到。\(^\text{[254]}\)  

短周期彗星（short-period comets）的轨道周期小于两百年；长周期彗星（long-period comets）的轨道周期则可长达数千年。天文学家普遍认为：短周期彗星起源于柯伊伯带（Kuiper belt），而像海尔–波普彗星（Hale–Bopp）这样的长周期彗星，起源于奥尔特云（Oort cloud）。许多彗星群（如克罗伊茨掠日群彗星，Kreutz sungrazers）被认为源自同一母体的分裂。\(^\text{[255]}\)部分具有双曲线轨道（hyperbolic orbits）的彗星可能起源于太阳系之外，但由于轨道测定困难，其确切来源尚不确定。\(^\text{[256]}\)那些因长期受太阳加热、其内部挥发物几乎被完全驱逐的古老彗星，通常被归类为小行星（asteroids）。\(^\text{[257]}\)  
\subsection{边界区域与不确定性}
\begin{figure}[ht]
\centering
\includegraphics[width=8cm]{./figures/8825dddc6727ab2a.png}
\caption{处于星际介质中的太阳系（左侧），示意图以对数刻度显示了各区域及其距离。} \label{fig_Solar_21}
\end{figure}
太阳系的外层区域仍有大量未知。距离太阳 100 个天文单位（AU）以外的空间几乎尚未被探索，因此对该区域的研究极为困难。人类目前对这一区域的认识，主要依赖于那些轨道被扰动、偶然靠近太阳的天体所提供的推测。即便如此，探测这些天体往往也只能在它们变得足够明亮、被观测为彗星时才有可能实现。\(^\text{[258]}\)可以确定的是，太阳系外层仍存在大量尚未被发现的天体。\(^\text{[259]}\)  

据估计，太阳的引力势能范围在约 两光年（125,000 AU）内仍主导着周围恒星的引力影响。太阳的希尔球（Hill sphere）—— 即太阳的引力势与银河系的引力势相等的边界 —— 被认为正好包围着奥尔特云（Oort cloud）其范围可能延伸至距离太阳约 230,000 AU之远。\(^\text{[8][27][260]}\)  

日球层（heliosphere）与希尔球的边界，代表了太阳系与银河系环境之间的交界处。在那里，太阳的引力势、星际介质（interstellar medium）与银河系整体引力势达到平衡。  
\subsubsection{日球层的边缘}
\begin{figure}[ht]
\centering
\includegraphics[width=8cm]{./figures/14aa982d0cc2ac5f.png}
\caption{太阳磁层与日球层顶（heliosheath）的示意图} \label{fig_Solar_22}
\end{figure}
太阳的恒星风气泡——称为日球层（heliosphere）—— 是一个由太阳主导的空间区域，其边界位于终止激波（termination shock）处。根据太阳相对于局部静止标准（local standard of rest）的特有运动，这一边界在朝向星际介质的上风方向约距太阳 80–100 AU，而在下风方向约为 200 AU。\(^\text{[261]}\)在此处，太阳风（solar wind）与星际介质（interstellar medium）发生碰撞，其速度显著降低、密度增加并变得更加湍流，从而形成一个巨大的椭圆形结构 —— 称为日球层顶（heliosheath）。\(^\text{[261]}\)  

理论上，日球层顶的外形与行为极类似于彗星的尾部：在上风侧进一步延伸约 40 AU，而在下风侧则可拖曳至数千个天文单位之远。[263][264]来自“卡西尼号（Cassini）”与“星际边界探测器 （Interstellar Boundary Explorer）”的观测证据表明，日球层受到星际磁场（interstellar magnetic field）的约束作用，被压缩成一个泡状结构（bubble shape），[265][266]但其实际形状仍未确定。[267]  

日球层外缘的形态可能受到星际介质交互作用的流体动力学效应以及太阳南半球主导的磁场分布的共同影响。例如，它呈现出“钝形”，其中北半球的延伸范围比南半球多约 9 AU。[261]日球层顶（heliopause）被视为星际介质的起点。[93]在约 230 AU 处之外，存在所谓的弓形激波（bow shock）：这是太阳在穿越银河系时所留下的等离子体尾迹（plasma wake）。[268]位于日球层之外的大型天体依然受太阳引力束缚，但星际介质中的物质流动会使微小尺度天体的分布逐渐均匀化。[93]  
\subsection{天体邻域}
\begin{figure}[ht]
\centering
\includegraphics[width=6cm]{./figures/b0377718ca7ecd85.png}
\caption{局部星际云（Local Interstellar Cloud）、G 云（G-Cloud）及其周围恒星的示意图。截至 2022 年，太阳系在这些星际云中的确切位置仍是天文学中尚未解决的问题。[269]} \label{fig_Solar_23}
\end{figure}
在距离太阳 10 光年范围内，恒星相对稀少；其中最近的是半人马座 α（三合星系统，Alpha Centauri），距离约 4.4 光年，可能位于“局部泡”（Local Bubble）的 G 云（G-Cloud） 内。[270]半人马座 α A 与 α B 是一对类似太阳的紧密双星，而距离太阳最近的恒星 —— 红矮星比邻星（Proxima Centauri）—— 则以约 0.2 光年的距离环绕它们运行。2016 年，天文学家发现一颗潜在可居住的系外行星环绕比邻星运行，被命名为 **比邻星 b（Proxima Centauri b），这是目前距离太阳最近的已确认系外行星。[271]  

太阳系被局部星际云（Local Interstellar Cloud）包围，但尚不清楚它是嵌入于该星际云之中，还是仅位于其边缘之外。[272]在距离太阳 300 光年范围内，存在多个星际云——统称为局部泡（Local Bubble）。[272]这一结构呈“沙漏形空腔”或“超泡”，在星际介质中横跨约 300 光年。其内部充满高温等离子体，表明该结构可能由多次近期超新星爆发所形成。[273]  

与邻近的更大尺度结构相比，局部泡只是一个较小的超泡。在它的邻域中存在更广阔的结构：Radcliffe 波（Radcliffe Wave）与分裂线性结构（Split Linear Structure，原 Gould 带），二者的长度都达到数千光年。[274]这些结构共同属于猎户臂（Orion Arm）的一部分，该旋臂中包含了肉眼可见的银河系大部分恒星。[275]  

恒星通常在星团（star clusters）中形成，之后再逐渐解体为共动星群（co-moving associations）。其中一个肉眼可见的显著星群是大熊座移动星群（Ursa Major moving group），距离约 80 光年，位于局部泡内部。最近的星团是金牛座毕宿星团（Hyades），位于局部泡的边缘。距离最近的恒星形成区包括：南冕座分子云（Corona Australis Molecular Cloud）、蛇夫座ρ云复合体（Rho Ophiuchi cloud complex）以及金牛座分子云（Taurus molecular cloud）；其中金牛座分子云位于局部泡之外，是 Radcliffe 波的一部分。[276]  

恒星近掠事件（stellar flybys）——即恒星经过距离太阳 0.8 光年以内 —— 大约每 10 万年 发生一次。最近一次测得的事件是朔尔茨星（Scholz’s Star），约在7 万年前曾接近至5 万 AU，很可能穿过了奥尔特云（Oort cloud）的外层。[277]每10 亿年约有1\%的概率会有恒星掠过太阳100 AU范围内，这种事件可能扰乱整个太阳系的轨道结构。[278]  
\subsection{银河系位置}
\begin{figure}[ht]
\centering
\includegraphics[width=6cm]{./figures/aa4fe6a2c1af5602.png}
\caption{银河系示意图，标注了主要的星系结构特征以及太阳系的相对位置。} \label{fig_Solar_24}
\end{figure}
太阳系位于银河系（Milky Way）之中。银河系是一种棒旋星系（barred spiral galaxy），直径约10 万光年，包含超过1000亿颗恒星。[279]太阳系所在的区域属于银河系外侧的一个旋臂，称为猎户–天鹅臂（Orion–Cygnus Arm）或本地支臂（Local Spur）。[280][281]太阳属于薄盘恒星族（thin disk population）的一员，这类恒星绕银河盘面近距离运行。[282]  

太阳绕银河系中心运行的速度约为220 km/s，完成一周约需2.4 亿年。[279]这一周期被称为太阳系的银河年（galactic year）。[283]太阳在星际空间中的运行方向称为太阳顶点（solar apex），其位置位于武仙座（Hercules）方向，大致朝向目前明亮恒星织女星（Vega）所在的位置。[284]黄道面与银河盘面之间的夹角约为60°。[c]  

太阳以近圆形轨道绕银河中心（Galactic Center）运行，该中心处为人马座 A（Sagittarius A\) 超大质量黑洞所在。太阳与银河中心的距离约为26,660 光年，[286]其公转速度与银河旋臂的旋转速度大致相同。[287]如果太阳的轨道更靠近银河中心，附近恒星的引力扰动可能会扰乱奥尔特云（Oort Cloud）中的天体，从而使大量彗星被送入内太阳系，造成与地球的撞击，对生命带来潜在的灾难性影响。此外，银河中心的强烈辐射也可能干扰复杂生命的演化。[287]  

太阳系在银河系中的位置对地球生命的演化历史具有重要影响。银河旋臂中存在更高密度的超新星（supernovae）、引力不稳定区（gravitational instabilities）与高能辐射环境，这些因素都可能扰乱太阳系的稳定性。然而，由于地球位于本地支臂（Local Spur），而非频繁穿越旋臂的区域，这为地球生命的长期演化提供了相对稳定的条件。[287]不过，根据颇具争议的湿婆假说（Shiva hypothesis），太阳系相对于银河系其他部分位置的周期性变化可能解释地球上出现的周期性生物大灭绝事件。[288][289]  
\subsection{发现与探索}
人类对太阳系的认识在几个世纪中逐步积累与扩展。直到中世纪晚期至文艺复兴时期，从欧洲到印度的天文学家普遍认为 —— 地球静止于宇宙的中心，[290]并且与那些在天空中运动的神圣或以太天体本质上截然不同。尽管希腊哲学家萨摩斯的阿里斯塔克斯（Aristarchus of Samos）曾提出过关于太阳中心说（heliocentric）宇宙的推测，但尼古拉·哥白尼（Nicolaus Copernicus）是已知第一个建立可数学预测的日心体系模型的人。[291][292]  

日心说（Heliocentrism）并未立即战胜地心说（geocentrism），但哥白尼的工作得到了几位学者的支持，其中最著名的是约翰内斯·开普勒（Johannes Kepler）。开普勒在哥白尼模型的基础上进行了改进 —— 他允许行星轨道为椭圆形轨道，并结合了第谷·布拉赫（Tycho Brahe）的精确观测数据，编制了著名的《鲁道夫星表（Rudolphine Tables）》，这使得对当时已知行星位置的计算达到了前所未有的精确度。皮埃尔·伽森狄（Pierre Gassendi）于 1631 年利用该星表预测了水星凌日（transit of Mercury），杰里迈亚·霍罗克斯（Jeremiah Horrocks）则在 1639 年成功预测了金星凌日（transit of Venus）。这些成果为日心说与开普勒的椭圆轨道理论提供了强有力的验证。[293][294]  

在17 世纪，伽利略·伽利莱（Galileo Galilei）率先将望远镜引入天文学观测；他与西蒙·马里乌斯（Simon Marius）分别独立发现，木星（Jupiter）拥有四颗绕行的卫星。[295]克里斯蒂安·惠更斯（Christiaan Huygens）在此基础上发现了土星最大卫星——泰坦（Titan），并揭示了土星环的形状。[296]1677 年，爱德蒙·哈雷（Edmond Halley）
观测到水星凌日现象，由此意识到：过观测行星的太阳视差（solar parallax）—— 理想情况下为金星凌日—— 可以使用三角测量法精确地确定地球、金星与太阳之间的距离。[297]  
他的朋友艾萨克·牛顿（Isaac Newton）在 1687 年出版的巨著《自然哲学的数学原理（Principia Mathematica）》中，证明了天体并非与地球上的物体本质不同：相同的运动定律与引力法则同样适用于地球与天空中的天体。[59]: 142   
\begin{figure}[ht]
\centering
\includegraphics[width=6cm]{./figures/36a444aa023572a3.png}
\caption{} \label{fig_Solar_25}
\end{figure}
“太阳系（Solar System）”这一术语最早于1704 年进入英语语言体系，当时约翰·洛克（John Locke）用它来指代太阳、行星与彗星。[298]1705 年，爱德蒙·哈雷（Edmond Halley）意识到多次观测到的某颗彗星其实是同一个天体，它以75–76 年为周期规律性地回归。这是第一个证明除行星外，亦有天体周期性环绕太阳运行的证据，[299]尽管早在公元 1 世纪，塞内卡（Seneca）就曾在理论上推测过类似的观点。[300]1769 年对金星凌日（transit of Venus）的精确观测，使天文学家得以计算出平均地日距离（Earth–Sun distance）为93,726,900 英里（150,838,800 公里），仅比现代测定值高出 0.8\%。[301]  

天王星（Uranus）自 1690 年起偶有被观测到，甚至可能早于古代，但直到 1783 年才被确认为位于土星轨道之外的行星。[302]1838 年，弗里德里希·贝塞尔（Friedrich Bessel）成功测量了恒星视差（stellar parallax）—— 即由于地球绕太阳公转而导致的恒星位置的视在偏移，这为日心说提供了第一个直接且实验证据。[303]几年后的1846 年，海王星（Neptune）被确认为行星，其发现得益于它对天王星轨道的微弱但可测的引力扰动。[304]而对水星轨道异常（Mercury’s orbital anomaly）的观测曾促使天文学家假设存在比水星更靠内的行星——火神星（Vulcan），但这一假说最终在阿尔伯特·爱因斯坦（Albert Einstein）于1915 年提出的广义相对论（general relativity）中被否定。[305]  

在 20 世纪，人类开始了太阳系的太空探索（space exploration）。自 20 世纪 60 年代起，人类便将望远镜送入太空轨道。[306]到 1989 年，八大行星均已被探测器访问过。[307]探测器曾从彗星[308]与小行星[309]采集样本返回地球，亦曾穿越太阳日冕（corona）[310]，并访问过两颗矮行星：冥王星（Pluto）与谷神星（Ceres）。[311][312]为节省燃料，一些任务利用引力助推（gravity assist）机动，例如旅行者一号与二号（Voyager 1 & 2）在飞掠外行星时获得加速，[313]而帕克太阳探测器（Parker Solar Probe）则在飞越金星后减速，进入更靠近太阳的轨道。[314]  

人类在 20 世纪 60–70 年代通过阿波罗计划（Apollo program）成功登月，[315]并计划在 2020 年代通过阿耳忒弥斯计划（Artemis program）重返月球。[316] 20 至 21 世纪的发现推动了“行星”定义的重新修订（2006 年），导致冥王星被降级为矮行星（dwarf planet），[317]并激发了人类对海王星外天体（trans-Neptunian objects）的更深研究兴趣。[318]  
\begin{figure}[ht]
\centering
\includegraphics[width=8cm]{./figures/5e9a198bd9426927.png}
\caption{} \label{fig_Solar_26}
\end{figure}
\subsection{参见}
\begin{itemize}
\item 行星际航行（Interplanetary spaceflight） —— 载人或无人进行的恒星或行星间旅行  
\item 太阳系内受引力整形的天体列表（List of gravitationally rounded objects of the Solar System）  
\item 太阳系极值列表（List of Solar System extremes）  
\item 按大小排列的太阳系天体列表（List of Solar System objects by size） —— 太阳系中最大的天体  
\item 太阳系地质特征列表（Lists of geological features of the Solar System） ——  
  小行星、卫星及除地球以外行星的地质特征目录  
\item 太阳系纲要（Outline of the Solar System） —— 太阳系的总体概览与主题导引  
\item 行星助记语（Planetary mnemonic） —— 用于记忆太阳系行星顺序的短语  
\item 虚构作品中的太阳系（Solar System in fiction） —— 文学、影视及艺术作品中的太阳系描绘  
\end{itemize}

